%% ============================================================
\section{Testing Methodology}
\label{sec:methodology}

\subsection{Juliet Benchmark Protocol}

The Juliet benchmark operates in two modes:

\textbf{Full mode}: All 305 enabled rules are applied to all 118 CWEs.  This
measures overall tool behavior but introduces noise from rules
triggering on CWEs they were not designed to detect.

\textbf{Fast mode}: Per-CWE rule manifests restrict each scan to only
the CERT C rules mapped to that CWE.  A mapping script generates 147
TOML manifest files from the official CERT C rule-to-CWE mapping.
This eliminates cross-CWE noise and reduces scan time by approximately
10$\times$.  All results from v0.3.20 onward use fast mode.

For each CWE directory, the benchmark invokes SqC with cross-file
analysis enabled, then classifies each violation by file name:
violations in \texttt{*bad*} files are True Positives; violations in
\texttt{*good*} files are False Positives.  Results are stored in a
SQLite database with per-violation granularity.

\textbf{TP rate} is defined as $\frac{\text{TP}}{\text{TP}+\text{FP}}$
(precision), not recall.  A ``bad'' file may contain multiple violation
types, but only the intended CWE is scored as TP.

\subsection{Real-World Benchmark Protocol}

Seven open-source C projects (\cref{tab:realworld-codebases}) are
analyzed with SqC, cppcheck (v2.10), and clang-tidy (v21.1.6).

\begin{table}[t]
\centering
\footnotesize
\begin{tabular}{@{}lrrl@{}}
\toprule
\textbf{Project} & \textbf{C Files} & \textbf{LOC} & \textbf{Commit} \\
\midrule
libcrc & 9 & 1,034 & \texttt{7719e21} \\
lua & 33 & 31,637 & \texttt{40b76de2} \\
raylib & 17 & 56,107 & \texttt{962bbfc} \\
mosquitto & 120 & 39,368 & \texttt{d3ee5c5} \\
curl & 222 & 186,220 & \texttt{3e198f75} \\
sqlite & 125 & 218,733 & \texttt{b1a73ba34} \\
hostap & 430 & 589,724 & \texttt{dcee6043} \\
\midrule
\textbf{Total} & \textbf{956} & \textbf{1,122,823} & \\
\bottomrule
\end{tabular}
\caption{Real-world benchmark codebases.}
\label{tab:realworld-codebases}
\end{table}

Real-world analysis revealed robustness requirements absent from the
Juliet benchmark.  Industrial embedded codebases frequently contain
comments written in non-Latin scripts: running SqC against a battery
management system codebase with Chinese-language function comments
uncovered a latent crash in three rule checkers caused by unsafe
UTF-8 string slicing (\cref{sec:architecture}).  The benchmark suite
itself later exposed two more analyzer-crashing defect classes that
Juliet never exercises.  hostapd's configuration parser contains an
\texttt{else if} chain several thousand branches deep, which overflowed
the stack in every rule checker that walked the AST by self-recursion
(three checkers crashed before the recursive walks were converted to
explicit-stack iteration); and a curl unit test containing a string
literal with a \texttt{]} character preceding the \texttt{[} of an
array subscript triggered an out-of-bounds slice in a bounds-checking
rule.  Juliet's test cases are structurally simple ASCII programs, so
none of these defect classes---encoding, nesting depth, or adversarial
token ordering---are visible to benchmark-only validation; all of them
surfaced exclusively through real-world analysis.

These experiences motivated a methodological decision: the real-world
suite is run under the same protocol rigor as Juliet.  Every run is
tied to a version bump and commit SHA, results are stored in the same
benchmark database, and per-rule per-project deltas are reviewed
between versions.  This discipline pays for itself: version-to-version
real-world comparisons have repeatedly surfaced both crashes and
per-rule regressions that the Juliet suite---where the same rule
changes showed neutral or positive movement---did not detect.  In our
experience Juliet should be treated as a rough approximation of
real-world behavior, useful for regression-testing detection logic but
neither sufficient to validate analyzer robustness nor predictive of
real-world precision (\cref{sec:limitations}).

An important caveat for interpreting the cross-tool results: SqC
implements 311 rules while cppcheck and clang-tidy implement
approximately 20 each.  The $\sim$30$\times$ difference in raw
violation counts reflects coverage breadth, not false positive rate.

\subsection{Source Scope and SLOC Accounting}
\label{sec:sloc-scope}

Precision, recall, and finding-density figures are only meaningful when
the line-count denominator matches the set of files actually evaluated.
A real-world repository ships substantial code outside the analysis
target: vendored third-party dependencies, test harnesses, build
tooling, language bindings, generated amalgamations, and
platform-specific backends that are not compiled in the reference
(Linux) build configuration.  Counting source lines over the whole
repository tree---as a naive \texttt{cloc} of the checkout would---
inflates the denominator and understates the true defect density of the
code under analysis.

We therefore define, per project, an \emph{evaluated scope}: the set of
first-party C translation units and headers that (a)~are built in the
reference configuration and (b)~were adjudicated in the ground-truth
audit (\cref{sec:realworld-results}).  All real-world SLOC figures in
this paper are reported over exactly this set, counting non-blank,
non-comment source lines.  \Cref{tab:eval-scope} lists the scope, the
audited commit, and the resulting SLOC for each project, together with
the selection criterion applied.  The exclusions are deliberate and
recorded with each codebase: for SQLite we evaluate the \texttt{src/}
core plus shipped \texttt{ext/} extensions but exclude the Tcl test
harness, build tooling, and the WASM/JNI bindings; for curl we exclude
the fourteen Windows/macOS backend files (\texttt{schannel}, Apple
Secure Transport, Win32) that the Linux build never compiles; for
Mosquitto we evaluate the \texttt{lib/} client library and \texttt{src/}
broker but exclude the bundled dependencies, test client, and plugin
examples; for Lua we evaluate the whole interpreter tree but exclude
the \texttt{ltests.c}/\texttt{ltests.h} test-harness pair; for raylib
we evaluate the core module set plus all platform backends (built
under \texttt{PLATFORM\_DESKTOP}, \texttt{ANDROID}, \texttt{DRM}, and
\texttt{WEB} configurations combined) but exclude the vendored
third-party decode/GL-loader headers; for hostap we evaluate the
shipped \texttt{src/}, \texttt{wpa\_supplicant/}, and \texttt{hostapd/}
trees (the station and access-point daemons and their shared library)
but exclude \texttt{tests/}, the separate \texttt{wlantest/} monitoring
tool, and the example/bindings directories that ship with the
repository but not with either daemon; for seL4 we evaluate the
\texttt{src/} kernel tree only, excluding \texttt{libsel4/} (the
userspace binding library, not kernel code), \texttt{tools/}, and
\texttt{manual/}.  libcrc is small enough that its entire shipped,
tooling, and test tree is in scope.  pure-ftpd is scoped differently
still: rather than a whole-project audit, it is our SQL-client-API
oracle (\cref{sec:ground-truth-snapshot}), scoped to exactly the two
files that motivated onboarding it---\texttt{src/log\_mysql.c} and
\texttt{src/log\_pgsql.c} plus their four headers, the MySQL/PostgreSQL
authentication backends that call \texttt{mysql\_real\_query}/
\texttt{PQexec} as a genuine SQL client---not the whole daemon, which
is scanned and tracked in the benchmark suite but not yet labeled.
Across the nine audited projects, the evaluated scope is substantially
smaller than the corresponding whole-repository line count; for
Mosquitto it is smaller by nearly a factor of three, so the choice of
denominator materially changes any per-KLOC metric.

\begin{table}[t]
\centering
\footnotesize
\begin{tabular}{@{}llrrl@{}}
\toprule
\textbf{Project} & \textbf{Commit} & \textbf{Files} & \textbf{SLOC} & \textbf{Scope criterion} \\
\midrule
\multicolumn{5}{@{}l}{\emph{Audited (C89/C90-style)}} \\
libcrc    & \texttt{7719e21} &  19 &     751 & whole shipped + tooling + test tree \\
mosquitto & \texttt{d3ee5c5} & 142 &  24,704 & \texttt{lib/} + \texttt{src/}; excl.\ deps, test, plugins \\
curl      & \texttt{3e198f7} & 427 & 110,615 & \texttt{lib/} + \texttt{src/}; excl.\ 14 Win/macOS files \\
sqlite    & \texttt{b1a73ba} & 220 & 161,312 & \texttt{src/} core + \texttt{ext/}; excl.\ test/tooling \\
hostap    & \texttt{dcee6043} & 736 & 442,335 & \texttt{src/}+\texttt{wpa\_supplicant/}+\texttt{hostapd/}; excl.\ tests/wlantest/examples \\
\cmidrule(l){1-5}
\textbf{Subtotal (C89/C90)} &        & \textbf{1,544} & \textbf{739,717} & \\
\midrule
\multicolumn{5}{@{}l}{\emph{Audited (C99 idiom corpora)}} \\
lua       & \texttt{40b76de2} &  66 &  20,986 & whole interpreter tree; excl.\ \texttt{ltests.*} \\
raylib    & \texttt{962bbfc}  &  23 &  38,473 & core + 10 platform backends; excl.\ vendored decode/GL headers \\
\cmidrule(l){1-5}
\textbf{Subtotal (C99)} &        & \textbf{89} & \textbf{59,459} & \\
\midrule
\multicolumn{5}{@{}l}{\emph{Audited (freestanding kernel)}} \\
seL4      & \texttt{1326364}  & 183 &  37,560 & \texttt{src/} kernel only; excl.\ \texttt{libsel4/}, tools, manual \\
\midrule
\multicolumn{5}{@{}l}{\emph{Audited (scoped SQL-client oracle, not whole-project)}} \\
pure-ftpd & \texttt{cc28bff5} &   6 &   1,436 & \texttt{log\_mysql}/\texttt{log\_pgsql} + headers only; rest of daemon unlabeled \\
\midrule
\multicolumn{5}{@{}l}{\emph{Total, all 9 audited}} \\
\textbf{Total}    &        & \textbf{1,822} & \textbf{838,172} & \\
\midrule
\multicolumn{5}{@{}l}{\emph{Planned}} \\
ReactOS   & ---               & --- & --- & Windows module subset (TP/FP scope) \\
\bottomrule
\end{tabular}
\caption{Evaluated source scope per project: SLOC counted over the
audited file set, not the whole repository.  Lua and raylib are
fully adjudicated (23/23 files for raylib, run \#66); hostap is
fully adjudicated across all 736 in-scope files (337 batches); seL4
is fully scoped to its 183-file kernel tree; pure-ftpd is
intentionally scoped to its 6-file SQL-client oracle rather than the
whole daemon.  The Planned row is reserved for the not-yet-started
Windows target.}
\label{tab:eval-scope}
\end{table}

hostap, seL4, and pure-ftpd's evaluated scopes and adjudication status
are described above; hostap is a full-corpus audit (all 736 in-scope
files, recall as well as precision) like lua and raylib, while seL4
and pure-ftpd are both deliberately scoped precision-only samples
(\cref{sec:ground-truth-snapshot})---seL4's sample targets the
MSC12-C empty-body/no-effect-code class specifically rather than
every rule, and pure-ftpd's does not extend past its two SQL-client
files to the rest of the daemon.  A Windows-API corpus (ReactOS)
remains under consideration to exercise the \texttt{WIN}-category rules
(\cref{sec:limitations}).  Because ReactOS exceeds the entire current
suite by several multiples in line count, it is anticipated as a
precision (true/false-positive) target with a module-scoped audit rather
than a whole-tree recall audit.

\subsection{Ground-Truth Adjudication Is a Snapshot, Not an Asymptote}
\label{sec:ground-truth-snapshot}

The real-world precision and recall figures in this paper
(\cref{sec:realworld-results,sec:provenance-frontier}) come from a
ground-truth oracle: individual findings manually labeled TP or FP and
stored keyed on the exact tuple (project, commit, file, line, rule).
Three properties of how that oracle is built limit what its numbers
can be read to mean.

First, adjudication is finding-based and incremental, never an
exhaustive re-derivation of every possible defect in a file. An audit
labels the findings a specific tool version actually emitted; it does
not independently enumerate every true vulnerability the file
contains.

Second, false-negative hunting during an audit is scoped to whatever
bug categories that audit's brief called out, not to every rule sqc
implements. For example, the hostap audit brief targets resource
leaks, unchecked allocations, unclosed file descriptors, and
unvalidated wire-length fields; it does not target double-free
detection. Across 30 audit batches and a full source read-through of
\texttt{ieee802\_11.c}, zero MEM31-C false negatives were ever
recorded for that file---not because the file has none, but because
double-free was never in that read-through's scope. A rule's absence
from a file's recorded false negatives is evidence about audit scope,
not evidence the rule has no misses there.

Third, and as a consequence of the first two, reported precision and
recall should be read as ``measured against ground-truth snapshot $Y$
as adjudicated,'' not as an asymptotic property of the rule. When a
rule's detection logic changes, its new findings do not automatically
inherit labels from the old ones: they fall outside the labeled set
and are excluded from the precision/recall denominator---in either
direction---until re-adjudicated. This can make a metric look flat
while the unlabeled coverage underneath it silently narrows. Concretely,
during the v0.4.176 session that motivated this caveat, a MEM31-C
scoping fix roughly doubled the rule's raw finding count in the
labeled-sample subset (2,206 $\to$ 4,429) while the reported
precision/recall for that rule appeared unchanged, because none of the
$\sim$2,244 new findings had yet been adjudicated into the denominator.
Any historical progression-table row (\cref{fig:tp-rate-progression})
for a rule whose detection logic changed since its last full
adjudication pass should therefore be treated as provisional pending
delta-adjudication, not as a confirmed measurement at the newer
version.
